\documentclass[conference]{IEEEtran}
\IEEEoverridecommandlockouts
% The preceding line is only needed to identify funding in the first footnote. If that is unneeded, please comment it out.
\usepackage{cite}
\usepackage{amsmath,amssymb,amsfonts}
\usepackage{algorithmic}
\usepackage{graphicx}
\usepackage{subcaption}
\usepackage{textcomp}
\usepackage{xcolor}
\usepackage[backend=biber, style=ieee]{biblatex}
\addbibresource{references.bib}
\usepackage[utf8]{inputenc} % Para que las tildes y ñ funcionen
\usepackage[T1]{fontenc}    % Mejor codificación de fuente
\usepackage[spanish]{babel}
\usepackage{float}
\def\BibTeX{{\rm B\kern-.05em{\sc i\kern-.025em b}\kern-.08em
    T\kern-.1667em\lower.7ex\hbox{E}\kern-.125emX}}
\begin{document}

\title{Análisis de Imágenes Satelitales con la Transformada de Fourier\\}

\author{\IEEEauthorblockN{1\textsuperscript{ro} Argueta Flores, Alison Aracely}
\IEEEauthorblockA{\textit{Ingeniería Informática} \\
\textit{Universidad Centroamericana (UCA)}\\
Antiguo Cuscatlán, El Salvador \\
00076422@uca.edu.sv}
\and
\IEEEauthorblockN{2\textsuperscript{do} Chávez Grande, Claudia María}
\IEEEauthorblockA{\textit{Ingeniería Informática} \\
\textit{Universidad Centroamericana (UCA)}\\
Antiguo Cuscatlán, El Salvador \\
00037221@uca.edu.sv}
\and
\IEEEauthorblockN{3\textsuperscript{ro} Hernández Cortez, Erika Sofía}
\IEEEauthorblockA{\textit{Ingeniería Informática} \\
\textit{Universidad Centroamericana (UCA)}\\
Antiguo Cuscatlán, El Salvador \\
00161220@uca.edu.sv}
\and
\IEEEauthorblockN{4\textsuperscript{to} Martínez Avila, María Auxiliadora}
\IEEEauthorblockA{\textit{Ingeniería Informática} \\
\textit{Universidad Centroamericana (UCA)}\\
Antiguo Cuscatlán, El Salvador \\
00090121@uca.edu.sv}}

\maketitle

\begin{abstract}
El trabajo experimental realizado sobre el Movimiento Armónico Simple (MAS), se enfocó en dos sistemas fundamentales: el masa-resorte y el péndulo simple. El estudio tuvo como objetivos principales examinar cómo el cambio en la masa influye en el período de oscilación de un resorte y determinar de forma experimental la aceleración de la gravedad local.
La metodología se dividió en dos fases: primero, la determinación de la constante elástica ($k$) de un resorte mediante la Ley de Hooke para analizar la relación potencial entre la masa ($m$) y el período  ($T$). Segundo, se midieron los tiempos de oscilación de un péndulo simple variando su longitud ($L$) a ángulos pequeños ($\leq 5^{\circ}$). Los datos fueron procesados mediante regresiones lineales y gráficas de $T^2$ vs $L$ para obtener el valor de la gravedad ($g$).
Los resultados obtenidos se validaros comparando la gravedad calculada con el valor de referencia de ($9.78~m/s^2$). Al evaluar los errores porcentuales, se verifició que fueran menores al 10% para confirmar la precisión del experimento. El análisis permitió concluir que los datos experimentales respaldan satisfactoriamente los modelos teóricos de oscilaciones estudiados. 
\end{abstract}
\vspace{0.5em}
\begin{IEEEkeywords}
Movimiento Armónico Simple, constante elástica, péndulo simple, gravedad, oscilaciones.
\end{IEEEkeywords}

\section{Introducción}
El surgimiento de la observación de la Tierra se remonta y está ligada al avance y desarrollo de la fotografía, que sucedió a principios del siglo XIX. Gracias a los grandes avances tecnológicos ocurridos desde entonces, la capacidad para observar la tierra desde la orbitra terrestre baja y a través de satélites geoestacionarios ha progresado de manera extraordinaria y continua con resoluciones espaciales de mejor calidad, radiométricas y espectrales, y tiempos de revisita más cortos \cite{emery2017}.

Años más tarde, en Estados Unidos se llevaron a cabo estudios sobre cómo los satélites en órbita terrestre podrían apoyar a la comunidad de pronóstico meteorológico para monitorar la Tierra; estos estudios permitieron que se concretarán proyectos para crear nuevos satélites que monitorearan la atmósfera. Es así como finalmente el primer TIROS (Television and Infrared Observation Satellite) fue lanzado y comenzó a operar en abril de 1960. El satélite era estabilizado por giro, lo que significaba que el sensor orientado a la Tierra solo observaba una parte limitada de la latitud terrestre. A su vez, la cámara fue diseñada para mostrar la capa de nubes de la Tierra como una especie de indicador de los sistemas meteorológicos sobre el planeta, es por eso que los primeros diseños de satélites estaban impulsados principalmente por consideraciones meteorológicas y la necesidad de mejorar los pronósticos \cite{emery2017}.

Puesto que todas las imágenes del satélite TIROS eran analógicas, no se logró efectuar las correcciones de las distorsiones geométricas mediante métodos digitales, y el mapeo se realizaba superponiendo "cuadrículas deformadas" que se ajustaran mejor a la orientación de los rasgos globales. Es por eso que la visualización de estas primeras imágenes de televisión satelital se convirtieron en un método de descubrimiento, encontrando imágenes que mostraban diversos procesos atmosféricos importantes, y a su vez, llegaron a ser el estándar para estudiar las condiciones de nubes atmosféricas y poder deducir sus conexiones con patrones climáticos globales \cite{emery2017}.

Como mencionado anteriormente, habían muchas limitantes entorno a los métodos analógicos pero eso cambió con el lanzamiento del primer Sistema de Escáner Multiespectral (MSS) del Landsat, en 1972, que contenía 4 bandas espectrales de unos 100 nm de ancho cada una, un tamaño de píxel de 80 m, y además, proporcionaba datos de imagen satelital de propósito general directamente en formato digital. Con esto inició la era moderna de la teledetección terrestre desde el espacio, siendo los sistemas de teledetección actuales diversos y con un rango de prestaciones muy variados \cite{schowengerdt2006}. Con el inicio de la era moderna de la teledetección terrestre empezaron a surgir nuevas problemáticas. 

El concepto de una imagen puede definirse como una función bidimensional, f(x, y), donde x e y son coordenadas espaciales (planas), y la amplitud de f en cualquier par de coordenadas (x, y) se denomina intensidad o nivel de gris de la imagen en ese punto, es cuando x, y y los valores de intensidad de f son todos cantidades finitas y discretas, que se llama a la imagen una imagen digital. Una imagen digital está compuesta por un número finito de elementos, cada uno de los cuales tiene una ubicación y un valor particulares. Estos elementos se denominan elementos de imagen, elementos pictóricos, pels y píxeles, siendo píxel el término más utilizado para designar los elementos de una imagen digital. Para crear una imagen digital, primero se convierten los datos continuos captados en un formato digital, lo cual requiere de dos procesos: muestreo y cuantificación, al pasar por la cuantificación, la precisión alcanzada depende en gran medida del contenido de ruido de la señal muestreada. Es aquí donde se presenta una de las grandes problemáticas en las imágenes digitales \cite{gonzalez2018}.

Una de las soluciones para la problemática mencionada anteriormente parte del Análisis Numérico, partiendo especificamente de Fourier. La solución que se plantea es examinar y purificar las señales digitales mediante su traslado al dominio de la frecuencia. A través de la Transformada Discreta de Fourier Bidimensional (DFT 2D), es posible descomponer la matriz espacial de la imagen satelital en sus componentes frecuenciales periódicas, aislando de forma selectiva el ruido y las texturas de alta frecuencia de las tendencias globales de baja frecuencia. Sin embargo, la computación directa de la DFT 2D presenta una complejidad algorítmica de $\mathcal{O}(M^2 N^2)$ para una imagen de dimensiones $M \times N$, siendo una carga aritmética que resulta restrictiva para los sistemas de procesamiento meteorológico en tiempo real \cite{gonzalez2018}. Es acá en donde entra la Transformada Rápida de Fourier (FFT), la cual es un algoritmo eficiente que es utilizado para calcular la DFT de una secuencia. Fue descrita por primera vez en el clásico artículo de Cooley y Tukey en 1965, siendo una idea originada del trabajo no publicado de Gauss en 1805. Puede ser descrito como un algoritmo de tipo ``divide y vencerás'' que descompone recursivamente la DFT en DFT más pequeñas para reducir el número de cálculos. Por lo que como resultado, logra reducir la complejidad de la DFT de \(O(n^2)\) a \(O(n \log n)\), siendo \(n\) el tamaño de los datos. Esta reducción en el tiempo de cálculo es significativa, particularmente para datos con un \(N\) grande, y por esta razón la FFT es eficiente para utilizar en imágenes satelitales \cite{kong2020python}.

En base a las necesidades de procesamiento descritas y bajo el enfoque del Análisis Numérico, el presente proyecto tiene como objetivo el diseño, implementación y evaluación de una solución computacional especializada en el filtrado en frecuencia de imágenes satelitales., utilizando el lenguaje de programación Python y los módulos optimizados de la biblioteca \texttt{numpy.fft}, el equipo implementó un algoritmo basado en un Filtro Pasa-Altas Ideal (IHPF) para reducir las componentes de baja frecuencia asociadas a fondos homogéneos y resaltar las estructuras de alta frecuencia, tales como frentes nubosos y contornos geográficos. La solución se integró en una interfaz interactiva con \texttt{ipywidgets}, que permite variar el radio de corte y observar la respuesta del sistema en tiempo real. Asimismo, se implementó una animación de reconstrucción progresiva y un análisis cuantitativo mediante el Teorema de Parseval para evaluar la conservación de la energía espectral, contrastando la eficiencia del algoritmo con sus limitaciones inherentes, como el efecto de Gibbs.


\section{Marco Teórico}

\subsection{Definición Matemática del Problema a Resolver}
\input{definicionmatematica}

\subsection{Representación Digital de Imágenes}

Una imagen digital en escala de grises puede modelarse matemáticamente como una
función discreta $I: \mathbb{Z}^2 \rightarrow \mathbb{Z}$, donde cada par de
índices enteros $(x, y)$ identifica la posición de un píxel en una cuadrícula
bidimensional de dimensiones $M \times N$, y el valor $I(x, y)$ representa la
intensidad de luminancia en dicho punto \cite{gonzalez2018}. En imágenes de 8
bits por canal, este valor pertenece al rango discreto $[0, 255]$, donde $0$
corresponde al negro absoluto y $255$ al blanco.

Una imagen a color no es una única función, sino la superposición de tres funciones independientes, cada una asociada a un canal de color; rojo, verde y azul (modelo RGB). En lugar de registrar un único valor de intensidad, cada píxel almacena tres valores distintos (uno para cada canal). La combinación de estos tres datos es lo que determina el color percibido en ese punto específico \cite{gonzalez2018}.

En la práctica, las imágenes satelitales son capturadas originalmente con
información de diversos canales que codifican distintas longitudes de onda del
espectro electromagnético. Para el procesamiento mediante la Transformada de
Fourier es necesario reducir esta representación de múltiples canales a una
sola capa de intensidad. Dicha conversión a escala de grises se realiza
mediante una combinación lineal ponderada de los canales rojo ($R$), verde
($G$) y azul ($B$), siguiendo los coeficientes establecidos por la
recomendación Rec.~601 de la \textit{International Telecommunication Union}
(ITU). Esta recomendación es un estándar técnico internacional, originalmente
diseñado para la televisión digital, que define con precisión cuánto debe
pesar cada canal de color al fusionarse en un solo valor de luminancia, de
manera que el resultado sea perceptualmente (proceso mediante el cual el cerebro interpretan información de su entorno a través de los sentidos \cite{concepto_percepcion}) equivalente al brillo que percibe el ojo humano \cite{gonzalez2018}:

\begin{equation}
    I_{\text{gris}}(x,y) = 0.299\, R(x,y) + 0.587\, G(x,y) + 0.114\, B(x,y)
    \label{eq:grises}
\end{equation}

Es importante aclarar que los coeficientes $0.299$, $0.587$ y $0.114$ no representan bits, ni profundidad de color. Simplemente, son pesos de mezcla, números entre 0 y 1 que indican qué proporción de cada canal contribuye al resultado final, y que suman exactamente $1.0$. Su origen es perceptual, no digital, reflejan que el ojo humano es
más sensible a la luz verde, moderadamente sensible al rojo, y mucho menos
sensible al azul \cite{gonzalez2018}. La imagen resultante
$I_{\text{gris}}(x,y)$ constituye el proceso de análisis espectral
tratado en las secciones siguientes.

% ─────────────────────────────────────────
\subsection{La Transformada Discreta de Fourier Bidimensional (DFT 2D)}

La Transformada Discreta de Fourier (DFT) es el instrumento central
del presente trabajo. Antes de tratar las imágenes, presenta su la versión más simple de ellas. Es decir, transforma una secuencia finita de $N$ muestras $x[n]$, en un nuevo
conjunto de $N$ coeficientes complejos $X[k]$, donde cada coeficiente está asociado a una frecuencia específica \cite{kulkarni, tan2020}:

\begin{equation}
    X[k] = \sum_{n=0}^{N-1} x[n]\, e^{-j2\pi kn/N}, \quad k = 0, 1, \ldots, N-1
    \label{eq:dft1d}
\end{equation}

Cuando el dato de entrada no es una secuencia de una sola dimensión (tiempo),
sino una imagen que se organiza en dos dimensiones espaciales (filas y columnas), la transformada debe generalizarse para operar sobre ambos índices simultáneamente. A esta generalización se le denomina \textit{caso
bidimensional} o \textit{DFT 2D}. El término ``bidimensional'' no se refiere a
una propiedad física de la imagen, sino a que el dominio de la función de
entrada tiene dos índices independientes, $x$ y $y$, en lugar de uno solo.
Dada una imagen discreta $I(x, y)$ de tamaño $M \times N$, su DFT 2D se define
como \cite{gonzalez2018, kulkarni}:

\begin{equation}
    F(u, v) = \sum_{x=0}^{M-1} \sum_{y=0}^{N-1} I(x,y)\,
              e^{-j2\pi\left(\frac{ux}{M} + \frac{vy}{N}\right)}
    \label{eq:dft2d}
\end{equation}

donde $u \in \{0, \ldots, M-1\}$ y $v \in \{0, \ldots, N-1\}$ son las
frecuencias espaciales en las direcciones horizontal y vertical,
respectivamente. El resultado $F(u, v)$ es una matriz de números complejos, de
la misma dimensión que la imagen original. Cada valor $F(u,v)$ no corresponde a
un solo píxel de la imagen original, sino que resume (mediante la suma de todos los píxeles $(x,y)$) qué tan presente está ese par particular de frecuencias $(u,v)$ en la imagen completa. En otras palabras, cada componente del espectro mide la fuerza con la que un patrón ondulatorio específico aparece repetido a lo largo de toda la imagen.

\subsubsection{Espectro de Magnitud y Fase}

Cada coeficiente $F(u, v)$ puede descomponerse en su parte real e imaginaria, o en su magnitud:

\begin{equation}
    |F(u,v)| = \sqrt{\operatorname{Re}(F(u,v))^2 + \operatorname{Im}(F(u,v))^2}
    \label{eq:magnitud}
\end{equation}

y en su fase:
\begin{equation}
    \angle F(u,v) = \arctan\!\left(\frac{\operatorname{Im}(F(u,v))}
                                        {\operatorname{Re}(F(u,v))}\right)
    \label{eq:fase}
\end{equation}

El espectro de magnitud $|F(u,v)|$ indica cuánta energía aporta cada componente
de frecuencia a la imagen. Las frecuencias bajas (componentes cercanas al
origen del espectro) corresponden a variaciones lentas y suaves de intensidad,
como grandes regiones homogéneas de océano o nubes. Las frecuencias altas,
ubicadas en la periferia del espectro, codifican los cambios abruptos de
intensidad, es decir, los bordes, contornos y texturas finas presentes en la imagen satelital \cite{gonzalez2018}. La Figura~\ref{fig:dft_spectrum} ilustra
de forma esquemática esta relación entre una imagen y su espectro de magnitud.

\begin{figure}[H]
\centering
\begin{subfigure}{0.46\textwidth}
        \centering
        \includegraphics[width=0.8\linewidth]{images/ft_circle_fourier.png}
        \caption{Espectro de Fourier correspondiente a un objeto circular.}
    \end{subfigure}
    \hfill
    \begin{subfigure}{0.46\textwidth}
        \centering
        \includegraphics[width=0.8\linewidth]{images/ft_oblique_rectangle_fourier.png}
        \caption{Espectro de Fourier correspondiente a un rectángulo inclinado.}
    \end{subfigure}

    \caption{Ejemplos de espectros de magnitud obtenidos mediante la Transformada de Fourier. La distribución de energía en el dominio de frecuencias depende de la geometría y orientación de las estructuras presentes en la imagen espacial. Adaptado de \cite{mathworks_fft}.}
    \label{fig:dft_spectrum}
\end{figure}

\subsubsection{Escalado Logarítmico del Espectro}

Dado que los coeficientes de Fourier pueden abarcar varios órdenes de magnitud,
la visualización directa de $|F(u,v)|$ no otorga la información suficiente. Los
componentes de baja frecuencia dominan visualmente y los de alta frecuencia
quedan invisibles. Para corregirlo, se aplica una transformación
logarítmica:

\begin{equation}
    S(u,v) = 20 \cdot \log\!\left(|F(u,v)| + 1\right)
    \label{eq:log}
\end{equation}

El término $+1$ evita la indeterminación del logaritmo cuando $|F(u,v)| = 0$.
Esta representación permite apreciar tanto las componentes dominantes como las
de menor energía \cite{pythonnumericalmethods}.

% ─────────────────────────────────────────
\subsection{La Transformada Rápida de Fourier (FFT)}

\subsubsection{Justificación mediante notación de complejidad asintótica}

En el Análisis de Algoritmos para describir cómo crece el número de operaciones necesarias de un
algoritmo a medida que aumenta el tamaño de su entrada, ignorando constantes y
términos de menor orden. Por ejemplo, un algoritmo $\mathcal{O}(N)$ duplica su
trabajo si se duplica $N$; uno $\mathcal{O}(N^2)$ lo cuadruplica
\cite{cooley1967}.

El cálculo directo de la DFT, según la ecuación~\eqref{eq:dft2d}, requiere
evaluar la doble suma para cada uno de los $M \times N$ puntos del espectro de
salida. Cada uno de esos puntos exige, a su vez, recorrer los $M \times N$
píxeles de entrada. Esto da como resultado un total de
$\mathcal{O}(N^4)$ operaciones para una imagen cuadrada de $N \times N$ píxeles.
Esta cantidad de operaciones crece rápidamente al duplicar
la resolución de la imagen, el tiempo de cómputo se multiplica por
aproximadamente dieciséis. Por ello se dice que el cálculo directo es computacionalmente prohibitivo, el tiempo requerido se vuelve irrealizable (horas o días) incluso para imágenes de resolución media \cite{cooley1967}.

\subsubsection{El algoritmo de Cooley-Tukey}

La Transformada Rápida de Fourier (FFT) es un algoritmo que calcula los mismos coeficientes que la DFT, pero aprovechando la simetría y
periodicidad de los factores exponenciales complejos $e^{-j2\pi kn/N}$ para
reducir el número de operaciones necesarias. El algoritmo
desarrollado por Cooley y Tukey en 1965 se basa en el paradigma de divide y conquista, en lugar de calcular una DFT de longitud $N$, la
descompone recursivamente en dos DFTs de longitud $N/2$ (una sobre las muestras de índice par y otra sobre las de índice impar) y combina ambos resultados parciales mediante un número reducido de sumas y multiplicaciones, en una
operación denominada mariposa (butterfly), debido a la forma
que adopta su diagrama de flujo de datos. Este proceso se repite recursivamente hasta llegar a subproblemas triviales de tamaño 1
\cite{cooley1967}.

Gracias a esta estrategia, el número de multiplicaciones necesarias se reduce
de $\mathcal{O}(N^2)$ a $\mathcal{O}(N \log N)$. Con el fin de observar la magnitud de esta mejora, se puede considerar un caso
para $N = 1024$ muestras, la DFT directa requiere del orden de un
millón de operaciones ($N^2 \approx 1.05 \times 10^6$), mientras que la FFT requiere del orden de diez mil ($N \log_2 N \approx 1.02 \times 10^4$); es
decir, la ganancia en velocidad es de aproximadamente cien veces
\cite{cooley1967}. La
Figura~\ref{fig:butterfly} ilustra el diagrama de mariposa característico de este algoritmo.

\begin{figure}[H]
\centering
\includegraphics[width=0.3\textwidth]{images/Butterfly-FFT.png}
\caption{Ejemplo ilustrativo de una imagen y su espectro de magnitud en el dominio de Fourier. Las componentes de baja frecuencia se concentran alrededor del centro del espectro, mientras que las altas frecuencias aparecen alejadas de éste. Adaptado de \cite{mathworks_fft}.}
\label{fig:butterfly}
\end{figure}

\begin{figure}[H]
\centering
\includegraphics[width=0.3\textwidth]{images/DIT-FFT-butterfly.jpg}
\caption{Diagrama de mariposa (\textit{butterfly}) utilizado por el algoritmo FFT de Cooley-Tukey para combinar los resultados parciales obtenidos durante la descomposición recursiva de la DFT. Adaptado de \cite{wikipedia_butterfly}.}
\label{fig:/DIT-FFT-butterfly}
\end{figure}


\subsubsection{Extensión al caso bidimensional}

En el contexto de imágenes, la FFT 2D se implementa aplicando la FFT 1D sucesivamente a todas las filas y luego a todas las columnas de la imagen (o viceversa), gracias a una propiedad conocida como "separabilidad" de la
DFT 2D. Esto reduce la complejidad total de $\mathcal{O}(N^4)$ a
$\mathcal{O}(N^2 \log N)$, haciendo posible su aplicación práctica en
imágenes satelitales de alta resolución \cite{gonzalez2018}. El resultado
numérico es idéntico al de la DFT definida en la
ecuación~\eqref{eq:dft2d}; la única diferencia es la eficiencia
computacional con que se obtiene ese mismo resultado.

% ─────────────────────────────────────────
\subsection{Centrado del Espectro: Reordenamiento de Frecuencias}

La DFT sitúa la componente de frecuencia cero
$F(0,0)$ (la componente DC, que representa el valor medio de la imagen) en la
esquina superior izquierda del arreglo resultante. Las frecuencias aumentan hacia los bordes del arreglo y las frecuencias negativas aparecen en la mitad derecha e inferior \cite{pythonnumericalmethods}.

Esta distribución resulta poco intuitiva para el diseño de filtros en el dominio de frecuencia, ya que las frecuencias bajas y las altas quedan intercaladas en las esquinas del espectro. Por ello se aplica un reordenamiento que desplaza cíclicamente los cuadrantes del arreglo; es decir, cada
cuadrante se traslada a la posición opuesta, como si el arreglo se enrollara sobre sí mismo. De modo que la componente DC quede centrada en la posición $(M/2,\, N/2)$, con las frecuencias creciendo desde el centro hacia los bordes \cite{numpyfft}:

\begin{equation}
    F_c(u, v) = \operatorname{Centrado}\bigl(F(u,v)\bigr)
    \label{eq:fftshift}
\end{equation}

Este recentrado tiene dos ventajas para el proyecto. Primero, el
espectro centrado es visualmente interpretable, las frecuencias bajas aparecen en el centro y las altas en la periferia. Segundo, facilita el diseño de la máscara de filtrado, cuya geometría circular se define naturalmente respecto al centro del arreglo.

Antes de aplicar la transformada inversa, es indispensable deshacer este reordenamiento mediante la operación de centrado inversa, que restaura la distribución original de frecuencias esperada por el algoritmo de reconstrucción. Omitir este paso introduciría un desplazamiento de fase incorrecto en la imagen reconstruida \cite{numpyfft}.

% ─────────────────────────────────────────
\subsection{Filtrado en el Dominio de Frecuencias}

\subsubsection{Teorema de Convolución}

El fundamento teórico que justifica el filtrado en el dominio de frecuencias es
el \textit{teorema de convolución}. Este establece que la convolución de dos señales en el dominio espacial es equivalente a la multiplicación de sus
transformadas de Fourier en el dominio de frecuencias \cite{oppenheim2010}:

\begin{equation}
    I(x,y) * h(x,y) \;\longleftrightarrow\; F(u,v) \cdot H(u,v)
    \label{eq:convolucion}
\end{equation}

donde $h(x,y)$ es el filtro espacial y $H(u,v)$ es su respuesta en frecuencia. Aplicar un filtro en el dominio espacial mediante convolución requiere $\mathcal{O}(N^4)$ operaciones; hacerlo en el dominio de frecuencias mediante multiplicación elemento a elemento, aprovechando la FFT, lo reduce a
$\mathcal{O}(N^2 \log N)$ (enfoque adoptado en el proyecto) \cite{gonzalez2018}.

\subsubsection{Filtro Pasa-Altas Ideal (IHPF)}

Un filtro pasa-altas tiene por objetivo suprimir las componentes de baja frecuencia del espectro y conservar únicamente las de alta frecuencia. El filtro pasa-altas ideal (\textit{Ideal High-Pass Filter}, IHPF) se establece mediante una máscara binaria $H(u,v)$ basada en la \textit{distancia euclidiana}, es decir, la distancia en línea recta, medida entre cada punto del espectro $(u,v)$ y el centro del mismo $(u_0, v_0)$ \cite{gonzalez2018}:

\begin{equation}
    D(u,v) = \sqrt{(u - u_0)^2 + (v - v_0)^2}
    \label{eq:distancia}
\end{equation}

\begin{equation}
    H(u,v) = \begin{cases}
        0 & \text{si } D(u,v) < R \\
        1 & \text{si } D(u,v) \geq R
    \end{cases}
    \label{eq:ihpf}
\end{equation}

donde $R$ es el radio de corte. Todas las frecuencias dentro del disco de radio $R$ centrado en el origen espectral son anuladas; las que se encuentran fuera
de este disco son conservadas sin modificación. En la implementación, este radio se fija inicialmente en $R = 30$ píxeles, valor que
define el umbral entre las componentes consideradas de baja y alta frecuencia
para la imagen satelital procesada. La Figura~\ref{fig:ihpf_mask} muestra un
ejemplo esquemático de esta máscara circular.

\begin{figure}[H]
    \centering
    \includegraphics[width=0.5\textwidth]{images/ihpf_process.jpg}
    \caption{Aplicación del filtro pasa-altas ideal (IHPF) en el dominio de la frecuencia. De izquierda a derecha: (a) imagen satelital en escala de grises, (b) máscara circular del filtro pasa-altas ideal utilizada para eliminar las componentes de baja frecuencia mediante un radio de corte de $R=30$ píxeles, y (c) imagen reconstruida tras aplicar la transformada inversa de Fourier, donde se resaltan los detalles de alta frecuencia.}
    \label{fig:ihpf_mask}
\end{figure}

El espectro filtrado $G(u,v)$ se obtiene mediante la multiplicación elemento a
elemento entre el espectro centrado y la máscara:

\begin{equation}
    G(u,v) = F_c(u,v) \cdot H(u,v)
    \label{eq:filtrado}
\end{equation}

\subsubsection{Distinción respecto al Filtro Gaussiano}
\label{sec:ghpf}

Es importante distinguir el filtro implementado en este proyecto de un filtro
pasa-altas Gaussiano (\textit{Gaussian High-Pass Filter}, GHPF), cuya función
de transferencia es \cite{gonzalez2018}:

\begin{equation}
    H_G(u,v) = 1 - e^{-D(u,v)^2 / (2\sigma^2)}
    \label{eq:ghpf}
\end{equation}

A diferencia del IHPF, el filtro Gaussiano presenta una transición suave y
continua entre las frecuencias atenuadas y las conservadas, controlada por el
parámetro $\sigma$. El filtro ideal, al introducir una discontinuidad abrupta
en $D = R$, produce el fenómeno conocido como \textit{ringing} o efecto de Gibbs en la imagen reconstruida, artefactos oscilatorios visibles alrededor de los bordes. Esta es una limitante inherente al diseño del filtro empleado en
el presente trabajo.

\begin{figure}[H]
    \centering
    \includegraphics[width=0.5\textwidth]{images/Gibbs_phenomenon_50.jpg}
    \caption{Aproximación funcional de una onda cuadrada mediante una serie de Fourier con 25 armónicos. Se observa el fenómeno de Gibbs, caracterizado por las oscilaciones y el sobreimpulso que aparecen en las proximidades de las discontinuidades y que no desaparecen al incrementar el número de términos de la serie, sino que se concentran en una región cada vez más estrecha. Adaptado de \cite{wikipedia_gibbs}.}
    \label{fig:gibbs}
\end{figure}

% ─────────────────────────────────────────
\subsection{Transformada Inversa de Fourier Bidimensional (IFFT 2D)}

Una vez aplicado el filtro en el dominio de frecuencias, la imagen filtrada se
recupera en el dominio espacial mediante la Transformada Inversa de Fourier
Discreta (IDFT) 2D \cite{kulkarni, johnson}:

\begin{equation}
    I'(x,y) = \frac{1}{MN} \sum_{u=0}^{M-1} \sum_{v=0}^{N-1}
              G(u,v)\, e^{\,j2\pi\left(\frac{ux}{M} + \frac{vy}{N}\right)}
    \label{eq:ifft2d}
\end{equation}

donde $G(u,v)$ es el espectro filtrado definido en la
ecuación~\eqref{eq:filtrado}. La imagen reconstruida $I'(x,y)$ representa la
versión de la imagen original en la que las frecuencias bajas han sido
suprimidas, resaltando únicamente los bordes, contornos y texturas de alta
frecuencia.

En la implementación, primero se
revierte el reordenamiento de frecuencias descrito en la
ecuación~\eqref{eq:fftshift} para restaurar el ordenamiento estándar del arreglo filtrado, y luego se calcula la transformada inversa mediante el algoritmo equivalente a la FFT, aplicado en sentido inverso. Dado que la imagen
original es real y el filtro introducido es simétrico, los coeficientes imaginarios del resultado deberían ser nulos en teoría; en la práctica, errores de redondeo en aritmética de punto flotante producen residuos imaginarios de
magnitud despreciable, por lo que se toma el valor absoluto del resultado
\cite{numpyfft}:

\begin{equation}
    I'_{\text{real}}(x,y) = \left|\mathcal{F}^{-1}\!\left\{G(u,v)\right\}\right|
    \label{eq:abs}
\end{equation}

% ─────────────────────────────────────────
\subsection{Teorema de Parseval y Análisis de Energía}

El Teorema de Parseval establece que la energía total de una señal se conserva bajo la transformada de Fourier. La suma de los cuadrados de sus valores en el dominio espacial es igual a la suma de los cuadrados de las magnitudes de sus
coeficientes en el dominio de frecuencias \cite{ye2026, oppenheim2010}. Para el caso discreto bidimensional:

\begin{equation}
    \sum_{x=0}^{M-1}\sum_{y=0}^{N-1} |I(x,y)|^2
    = \frac{1}{MN}\sum_{u=0}^{M-1}\sum_{v=0}^{N-1} |F(u,v)|^2
    \label{eq:parseval}
\end{equation}

Una consecuencia de este teorema en el análisis del filtrado es que, al anular selectivamente los coeficientes de baja frecuencia mediante la máscara $H(u,v)$, disminuye la energía total del espectro; un efecto que, además, es posible cuantificar. La energía del espectro original y del espectro filtrado se
calculan como:

\begin{equation}
    E_{\text{original}} = \sum_{u}\sum_{v} |F_c(u,v)|^2
    \label{eq:energia_orig}
\end{equation}

\begin{equation}
    E_{\text{filtrada}} = \sum_{u}\sum_{v} |G(u,v)|^2
    \label{eq:energia_filt}
\end{equation}

El porcentaje de energía conservada tras el filtrado es:

\begin{equation}
    \eta = \frac{E_{\text{filtrada}}}{E_{\text{original}}} \times 100\%
    \label{eq:energia_conservada}
\end{equation}

En el proyecto, este análisis cuantitativo se recalcula para cada valor del
radio $R$ durante la animación, lo que permite observar cómo la energía
conservada decrece progresivamente a medida que se amplía la zona de exclusión.
Dado que las frecuencias bajas concentran la mayor parte de la energía de una
imagen natural, incluso un radio moderado de $R = 30$ píxeles puede eliminar un
porcentaje significativo de la energía total, conservando únicamente la
aportación estructural de alta frecuencia \cite{gonzalez2018}.

% ─────────────────────────────────────────
\subsection{Imágenes Satelitales GOES-East: Fuente de Datos}

Las imágenes procesadas en este proyecto provienen del satélite
\textit{GOES-East} (Geostationary Operational Environmental Satellite),
operado por la \textit{National Oceanic and Atmospheric Administration} (NOAA)
y distribuido públicamente a través de la plataforma de visualización
interactiva del Centro de Vuelos Espaciales George C. Marshall de la NASA, bajo
el programa \textit{Short-term Prediction Research and Transition} (SPoRT)
\cite{nasasport2026}.

Las imágenes fueron obtenidas del canal 7 del instrumento \textit{Advanced
Baseline Imager} (ABI), que opera en la banda del infrarrojo de onda corta a
$3.90\,\mu\text{m}$, con una resolución espacial nativa de 2 km por píxel
\cite{nasasport2026}. Este canal es especialmente adecuado para la detección de
nubes, niebla nocturna, focos de calor y estimación de vientos atmosféricos.

Desde la perspectiva del análisis de Fourier, las imágenes captadas por el canal 7 presentan características espectrales definidas. Grandes regiones de fondo homogéneo (océano, superficie terrestre) que concentran la energía en
frecuencias bajas, y estructuras bien delimitadas como bordes de sistemas
nubosos, costas y corrientes oceánicas superficiales, que aportan componentes
de alta frecuencia. Esta distribución espectral hace de este tipo de imágenes un caso adecuado para demostrar claramente el efecto del filtro pasa-altas implementado.

\section{Metodología}
\input{metodologia}

\section{Análisis y Resultados}
\input{resultados}


\subsection{Parte II. Péndulo simple}
\begin{table}[h!]
\centering
\caption{Datos para analizar la relación entre el período $T$ y la masa $M$ en el sistema masa-resorte [cite: 44].}
\label{tab:pendulo_compacta}
\resizebox{\columnwidth}{!}{% 
\begin{tabular}{|c|c|c|c|c|c|c|c|}
\hline
\textbf{$m \pm \Delta m$} & \textbf{$t_{1}$} & \textbf{$t_{2}$} & \textbf{$t_{3}$} & \textbf{$t_{4}$} & \textbf{$t_{5}$} & \textbf{$t \pm \Delta t$} & \textbf{$T \pm \Delta T$}  \\ 
\textbf{(g)} & \textbf{(s)} & \textbf{(s)} & \textbf{(s)} & \textbf{(s)} & \textbf{(s)} & \textbf{(s)} & \textbf{(s)}  \\ \hline
550 $\pm$ 0.1 & 6.13 & 6.17 & 6.73  & 6.65 & 6.09 & 6.35 $\pm$ 0.1 & 1.27 $\pm$ 0.31 \\ \hline
450 $\pm$ 0.1 & 6.13  & 6.15  & 5.79  & 5.90 & 5.73  & 5.94 $\pm$ 0.1  & 1.19 $\pm$ 0.19 \\ \hline
350 $\pm$ 0.1 &  4.97 & 5.09  & 4.95  & 5.01  & 4.94  & 4.99 $\pm$ 0.1  & 1.00 $\pm$ 0.061  \\ \hline
250 $\pm$ 0.1 & 4.40 & 4.32 & 4.07 & 4.17  & 4.45  & 4.28 $\pm$ 0.1 & 0.86 $\pm$ 0.16  \\ \hline
150 $\pm$ 0.1 & 3.07  & 3.48  & 3.09  & 3.28  & 3.38  & 3.26 $\pm$ 0.1  & 0.65 $\pm$ 0.18 \\ \hline

\end{tabular}
}
\end{table}

La tabla II presenta datos experimentales primarios, obtenidos durante el procedimiento. Estos resultados númericos recopilados serán utilizados para el análisis cuantitativo de la relación del período $T$ y la masa $M$. Se espera validar el modelo del oscilador armónico simple. Cada fila de la tabla corresponde a un valor de masa controlado con cinco mediciones del tiempo en el que esta oscila cinco veces. Se calcula entonces, un promedio de tiempo y se toma en cuenta una incerteza de 0.1 segundos. En base a estas mediciones, se calcula el período $T$ y su incerteza $\Delta T$ (incerteza obtenida con el cálculo de desviación estándar con una muestra). Las mediciones del tiempo fueron repetidas con la intención de reducir el impacto que los errores aleatorios tendrían en los resultados y para la obtención de una estimación estadística de la incerteza del período. La comparación del período con su incerteza estadística ($\frac{\Delta T}{T}$) muestra una incerteza relativa que varía en un rango de 6\% y 28\% según distintos períodos calculados. La dispersión entre datos indica que la precisión en las mediciones fue no uniforme durante el procedimiento experimental. Se  considera que esto puede ser atribuido a errores dominantes no instrumentales relacionados con la intervención humana durante la medición con el cronómetro, dado que estos procedimientos pueden introducir más riesgo de error que otros experimentos (como la utilización de una balanza estable, por ejemplo). No se ha sobreestimado la precisión de los datos reportados, considerando las cifras de los instrumentos utilizados hasta las operaciones realizadas a partir de los datos primarios.


Los hallazgos del procedimiento sugieren que el período aumenta con el incremento de la masa en el sistema. Lo evidenciado, entonces, concuerda cualitativamente con el modelo teórico conocido-propuesto. Inicialmente no se señalan comportamientos anómalos o datos atípicos en los resultados cuantitativos obtenidos en la tabla. Lo evidenciado en la tabla II permite analizar la dependencia.


\begin{table}[h!]
\centering
\caption{Datos para analizar la relación entre el período $T$ y la longitud $L$ en el péndulo simple[cite: 44].}
\label{tab:pendulo_compacta}
\resizebox{\columnwidth}{!}{% 
\begin{tabular}{|c|c|c|c|c|c|c|c|c|}
\hline
\textbf{$L \pm \Delta L$} & \textbf{$t_{1}$} & \textbf{$t_{2}$} & \textbf{$t_{3}$} & \textbf{$t_{4}$} & \textbf{$t_{5}$} & \textbf{$t \pm \Delta t$} & \textbf{$T \pm \Delta T$} & \textbf{$T^{2} \pm \Delta T^{2}$} \\ 
\textbf{(m)} & \textbf{(s)} & \textbf{(s)} & \textbf{(s)} & \textbf{(s)} & \textbf{(s)} & \textbf{(s)} & \textbf{(s)} & \textbf{(s$^2$)} \\ \hline
1.000 $\pm$ 0.001 & 10.05 & 10.12 & 9.98  & 10.03 & 10.07 & 10.05 $\pm$ 0.05 & 2.01 $\pm$ 0.01 & 4.04 $\pm$ 0.04 \\ \hline
0.900 $\pm$ 0.001 & 9.52  & 9.48  & 9.55  & 9.60  & 9.50  & 9.53 $\pm$ 0.04  & 1.91 $\pm$ 0.01 & 3.65 $\pm$ 0.04 \\ \hline
0.800 $\pm$ 0.001 & 8.95  & 9.02  & 8.88  & 8.92  & 8.98  & 8.95 $\pm$ 0.05  & 1.79 $\pm$ 0.01 & 3.20 $\pm$ 0.04 \\ \hline
0.700 $\pm$ 0.001 & 8.35  & 8.42  & 8.30  & 8.45  & 8.38  & 8.38 $\pm$ 0.06  & 1.68 $\pm$ 0.01 & 2.82 $\pm$ 0.03 \\ \hline
0.600 $\pm$ 0.001 & 7.75  & 7.80  & 7.72  & 7.85  & 7.78  & 7.78 $\pm$ 0.05  & 1.56 $\pm$ 0.01 & 2.43 $\pm$ 0.03 \\ \hline
0.500 $\pm$ 0.001 & 7.08  & 7.15  & 7.02  & 7.10  & 7.12  & 7.09 $\pm$ 0.05  & 1.42 $\pm$ 0.01 & 2.02 $\pm$ 0.03 \\ \hline
0.400 $\pm$ 0.001 & 6.32  & 6.38  & 6.28  & 6.35  & 6.30  & 6.33 $\pm$ 0.04  & 1.27 $\pm$ 0.01 & 1.61 $\pm$ 0.02 \\ \hline
\end{tabular}
}
\end{table}

\section{Conclusiones y Limitantes}

\subsection{Síntesis del trabajo realizado}

La Transformada
Rápida de Fourier constituye una herramienta eficaz para extraer información
estructural de imágenes satelitales que, de otro modo, permanece implícita en
el dominio espacial. A partir de una imagen del canal 7 del instrumento ABI a
bordo del satélite GOES-East, se logró transformar la información de
intensidad al dominio de frecuencias, diseñar e implementar un filtro
pasa-altas ideal sobre dicho espectro, reconstruir la imagen mediante la
transformada inversa, y cuantificar el efecto del filtrado mediante el Teorema
de Parseval. El resultado fue una imagen donde las regiones homogéneas de
fondo —características de superficies oceánicas y terrestres extensas— quedan
suprimidas, mientras que los bordes de formaciones nubosas, líneas de costa y
otras discontinuidades de intensidad emergen con mayor definición. Esto
confirma, en un caso de estudio concreto, la correspondencia teórica discutida
en el Marco Teórico entre frecuencias bajas y regiones de variación lenta, y
entre frecuencias altas y cambios abruptos de intensidad.

\subsection{Aportes del enfoque implementado}

Uno de los aportes más relevantes del proyecto fue la incorporación del
análisis energético basado en el Teorema de Parseval como complemento
cuantitativo al filtrado. Mientras que la inspección visual de la imagen
filtrada permite apreciar cualitativamente qué estructuras se resaltan, el
cálculo del porcentaje de energía conservada $\eta$ aporta una medida objetiva
de cuánta información se descarta en el proceso. Esta dualidad —resultado
visual y métrica numérica— enriquece la interpretación del experimento, ya que
permite distinguir entre un filtrado conservador, que retiene la mayor parte
de la energía y solo atenúa marginalmente las frecuencias bajas, y un
filtrado agresivo, que elimina una fracción sustancial de la energía total a
cambio de un realce más marcado de los bordes.

Asimismo, la implementación de la animación que varía el radio de corte $R$
de manera progresiva constituyó un aporte pedagógico importante del proyecto.
En lugar de evaluar el efecto del filtro en un único punto de operación, la
animación permite observar la transición continua entre la imagen original y
una versión donde prácticamente toda la información de baja frecuencia ha sido
eliminada. Esta visualización dinámica hace tangible un concepto que, en el
Marco Teórico, se presenta de forma estática: la relación inversa entre el
radio de corte y la energía conservada, así como el punto a partir del cual el
filtrado comienza a degradar significativamente la interpretabilidad de la
imagen.

\subsection{Ventajas observadas del filtro pasa-altas ideal}

La elección del filtro pasa-altas ideal, frente a alternativas como el filtro
Gaussiano discutido en la sección de Marco Teórico, presenta una ventaja
práctica relevante para los fines de este proyecto: su definición mediante una
máscara binaria simplifica tanto la implementación como la interpretación del
parámetro de control. El radio de corte $R$ tiene un significado geométrico
directo —el tamaño del disco de frecuencias eliminadas— que no requiere ajustar
un parámetro adicional de forma (como sí ocurre con la desviación estándar
$\sigma$ del filtro Gaussiano). Esta simplicidad resultó conveniente para el
experimento de barrido de $R$, en el que se requería un único parámetro libre
para generar la secuencia de imágenes filtradas.

\subsection{Limitantes técnicas y matemáticas}

A pesar de los resultados obtenidos, el proyecto presenta limitantes que deben
reconocerse con honestidad académica.

En primer lugar, la limitante más relevante desde el punto de vista teórico es
la presencia del efecto de Gibbs (\textit{ringing}), discutido en la sección
\ref{} del Marco Teórico. La discontinuidad abrupta de la máscara $H(u,v)$ en
$D(u,v) = R$ introduce artefactos oscilatorios alrededor de los bordes en la
imagen reconstruida, visibles como un patrón de anillos concéntricos. Este
fenómeno, inherente a cualquier filtro ideal y no exclusivo de esta
implementación, podría haberse mitigado adoptando un filtro con transición
suave, como el filtro Gaussiano o un filtro Butterworth, a costa de perder la
simplicidad conceptual del IHPF.

En segundo lugar, el radio de corte $R$ se definió y se exploró de forma
empírica (inicialmente fijado en $R = 30$ píxeles, y posteriormente recorrido
en un rango de $5$ a $150$ píxeles durante la animación), sin un criterio
matemático formal que determine un valor óptimo en función del contenido
espectral específico de cada imagen. Un valor de $R$ adecuado para resaltar
los bordes de un sistema nuboso extenso podría no ser igualmente adecuado para
una imagen con estructuras de menor escala espacial. Esta dependencia del
parámetro respecto al contenido de la imagen constituye una limitante para la
generalización del método a otros conjuntos de datos sin un proceso de
calibración previo.

En tercer lugar, la conversión de la imagen satelital original —que en su
forma nativa codifica información radiométrica de un canal infrarrojo
específico— a una escala de grises de 8 bits mediante los coeficientes de la
Rec.~601 introduce una pérdida de información. Dicha recomendación fue
diseñada originalmente para señales de video en el espectro visible (RGB), por
lo que su aplicación a un canal infrarrojo de onda corta, como el canal 7 del
instrumento ABI, no necesariamente preserva con fidelidad las relaciones
físicas de temperatura de brillo originales de la imagen satelital. El
análisis de Fourier desarrollado, por tanto, opera sobre una representación
visual derivada de los datos, y no sobre el dato radiométrico original sin
procesar.

En cuarto lugar, el filtro implementado actúa exclusivamente sobre la
componente de magnitud del espectro, dado que la máscara $H(u,v)$ se aplica de
manera idéntica a la parte real e imaginaria del espectro complejo. Esto
significa que, si bien el filtrado modifica qué frecuencias contribuyen a la
imagen reconstruida, no se realizó en este proyecto un análisis independiente
del comportamiento de la fase $\angle F(u,v)$, que también contiene
información estructural relevante, particularmente sobre la posición relativa
de los bordes. Una limitante metodológica del trabajo es, en consecuencia, no
haber explorado el rol de la fase de manera separada a la magnitud.

Finalmente, conviene señalar que la implementación se restringió al filtro
pasa-altas ideal con geometría circular, sin contrastar sus resultados frente
a otras geometrías de máscara (por ejemplo, máscaras rectangulares o
direccionales) que podrían resultar más adecuadas para resaltar estructuras
con una orientación predominante, como frentes nubosos alineados o líneas de
costa con una dirección preferente.

\subsection{Trabajo futuro}

A partir de las limitantes identificadas, se desprenden varias líneas de
mejora viables para una extensión futura de este trabajo. Una primera mejora
consistiría en sustituir el filtro pasa-altas ideal por su contraparte
Gaussiana o por un filtro Butterworth de orden ajustable, lo que permitiría
cuantificar experimentalmente, mediante el mismo análisis de energía basado en
el Teorema de Parseval, la reducción del efecto de Gibbs a cambio de una
transición más gradual entre frecuencias conservadas y eliminadas.

Una segunda mejora sería el desarrollo de un criterio automático de selección
del radio de corte $R$, basado, por ejemplo, en el porcentaje de energía que
se desea conservar ($\eta$ objetivo), de manera que el parámetro se ajuste de
forma adaptativa al contenido espectral de cada imagen en lugar de fijarse de
manera arbitraria.

Una tercera línea de trabajo futuro consistiría en extender el análisis a
series temporales de imágenes del mismo canal y la misma región geográfica,
capturadas en distintos instantes, lo que permitiría observar la evolución
del espectro de frecuencias asociada al desplazamiento de sistemas nubosos a
lo largo del tiempo, conectando este proyecto con la aplicación de detección
de estacionalidad en series de tiempo mencionada como alternativa en las
indicaciones del proyecto.

Por último, sería de interés incorporar al análisis la componente de fase del
espectro de Fourier, actualmente fuera del alcance de este trabajo, con el
fin de evaluar si la información posicional que ésta codifica aporta criterios
adicionales para la detección de bordes geográficos, complementando así el
análisis basado únicamente en la magnitud del espectro.

\section*{Normas de Citación y Referencias}

La plantilla numerará las citas de forma consecutiva entre corchetes [1]. La puntuación de la oración sigue al corchete [2]. Remítase simplemente al número de referencia, como en [3]—no utilice ``Ref. [3]'' ni ``referencia [3]'' excepto al principio de una oración: ``La referencia [3] fue la primera ...''.

Numere las notas al pie por separado mediante superíndices. Coloque la nota al pie real en la parte inferior de la columna en la que fue citada. No incluya notas al pie en el resumen ni en la lista de referencias. Utilice letras para las notas al pie de las tablas.

A menos que haya seis autores o más, indique los nombres de todos los autores; no utilice ``et al.''. Los artículos que no hayan sido publicados, incluso si han sido enviados para su publicación, deben citarse como ``inéditos'' [4]. Los artículos que hayan sido aceptados para su publicación deben citarse como ``en prensa'' [5]. Escriba con mayúscula solo la primera palabra del título de un artículo, excepto para nombres propios y símbolos de elementos.


\section{Referencias bibliográficas}

\printbibliography[title={Referencias}]


\end{document}
